% =====================================================================
% 5.1 Adversarial examples against vision models
%
% Source: ../aisb/5.1-adversarial-vision/section1_solution.py
% Figures (from each paper's arXiv source):
%   panda_577 / nematode_082 / gibbon_993   Goodfellow et al. 2015, Fig. 1
%   stop_graffiti_real, stop_sticker_attack Eykholt et al. 2018, Fig. 1
%   turtle_demo                             Athalye et al. 2018, Fig. 1
%   patch_banana_toaster                    Brown et al. 2017, Fig. 1
%   clip_apple_ipod                         Goh et al. 2021 (Distill), image
%                                           taken from The Guardian, 2021-03-08
% All percentages on these frames are quoted from those papers.
%
% NOTE: the lab does not implement FGSM or PGD. It runs Adam on the
% perturbation with an L2 penalty and an L-inf clamp. FGSM/PGD appear
% here as background only.
% =====================================================================

\section{5.1: Adversarial Examples (Vision)}

\begin{frame}
\frametitle{Advesarial Input for Image Classification}
\begin{center}
\setlength{\tabcolsep}{3pt}
\begin{tabular}{ccccc}
\includegraphics[width=0.24\textwidth]{panda_577} &
\raisebox{0.11\textwidth}{$+\;0.007\,\times$} &
\includegraphics[width=0.24\textwidth]{nematode_082} &
\raisebox{0.11\textwidth}{$=$} &
\includegraphics[width=0.24\textwidth]{gibbon_993} \\[2pt]
\small ``panda'' & & \small noise & & \small ``gibbon'' \\
\scriptsize 57.7\% confidence & & & & \scriptsize 99.3\% confidence
\end{tabular}
\figsrc{GoogLeNet on ImageNet. Figure from \textcite{goodfellow2015explaining}.}
\end{center}
\vspace{0.2cm}
\begin{itemize}
    \item We can compute the gradient
    of the loss with respect to pixels, and then perform gradient descent on the image.
    \item One can add extra terms to the loss to control the magnitude of the perturbation,
    such that the perturbation is invisible at first glance.
\end{itemize}
\end{frame}

\begin{frame}
\frametitle{It survives the physical world}
\begin{columns}[T]
\begin{column}{0.5\textwidth}
\centering
\includegraphics[height=3.0cm]{stop_graffiti_real}\hspace{2pt}%
\includegraphics[height=3.0cm]{stop_sticker_attack}
\figsrc{Left: real graffiti. Right: the attack.\\
\textcite{eykholt2018robust}.}
\end{column}
\begin{column}{0.5\textwidth}
\centering
\includegraphics[width=\linewidth]{turtle_demo}
\figsrc{3D-printed turtle. Red border: classified as rifle.\\
\textcite{athalye2018synthesizing}.}
\end{column}
\end{columns}
\vspace{0.3cm}
\begin{itemize}
    \item Black and white stickers on a real stop sign: read as Speed Limit 45
    in 15 of 15 lab photos, and 28 of 33 frames from a moving car.
    \item The turtle is classified as a rifle from almost every angle. The
    unperturbed print is classified as a turtle.
    \item Both were optimised over many camera angles and distances, not one
    photo. More on that shortly.
\end{itemize}
\end{frame}

\begin{frame}
\frametitle{It does not even need to be small}
\begin{center}
\includegraphics[width=0.58\textwidth]{patch_banana_toaster}
\figsrc{VGG16. Figure from \textcite{brown2017adversarial}.}
\end{center}
\begin{itemize}
    \item Once can learn \textbf{universal} adversarial example: a patch
    that can be added to the image to induce misclassification,
    rather than learning a per-image perturbation.
\end{itemize}
\end{frame}

\begin{frame}[plain]
\vfill
\centering
\makebox[\linewidth]{\includegraphics[width=\paperwidth]{cover_bright}}
\vfill
\end{frame}

\begin{frame}
\frametitle{It does not even need optimisation}
\begin{center}
\includegraphics[width=\textwidth]{clip_apple_ipod}
\figsrc{CLIP, zero-shot. Figure from \textcite{goh2021multimodal}, via The Guardian.}
\end{center}
\begin{itemize}
    \item Model here is \emph{CLIP}, which has trained on image--caption pairs.
    \item A side effect is that the vision model has learned to read text, providing
   another attack vector.
\end{itemize}
\end{frame}

\begin{frame}
\frametitle{Stated as an optimisation}
Freeze the weights. Let $\Loss(\xx, y)$ be the usual cross-entropy loss of
the classifier on image $\xx$ with label $y$.
\begin{block}{Targeted attack}
$$
\min_{\lVert\dd\rVert \,\le\, \eps}\; \Loss(\xx + \dd,\; y_\tgt)
$$
\end{block}
\begin{itemize}
    \item Perform constrained optimization on the perturbation $\dd$
    to minimize the loss $\Loss(\xx + \dd, y_\tgt)$, forcing the model
    to misclassify the perturbed image $\xx + \dd$ as $y_\tgt$.
    \item The norm bound $\eps$ is the attacker's stealth budget. 
    A lower $\eps$ means a more invisible perturbation, but makes the attack
    more difficult to achieve.
\end{itemize}
\end{frame}

% \begin{frame}
% \frametitle{Why these exist: dimension}
% \begin{itemize}
%     \item Consider a $224\times224$ RGB image = $150\,528$ inputs for perturbation. 
%     \item Many tiny nudges, each imperceptible, all add up to a large change.
% \end{itemize}
% \end{frame}

\begin{frame}
\frametitle{FGSM and PGD}
\begin{block}{FGSM: one step}
$$
\dd = -\,\eps\,\sign\!\big(\nabla_{\xx}\Loss(\xx, y_\tgt)\big)
$$
\end{block}
\begin{block}{PGD: many small steps, clip back into the box}
$$
\dd_{t+1} = \clip_{[-\eps,\,\eps]}\Big(
  \dd_t - \alpha\,\sign\big(\nabla_{\dd}\Loss(\xx + \dd_t, y_\tgt)\big)\Big)
$$
\end{block}
\begin{itemize}
    \item FGSM \parencite{goodfellow2015explaining} was used for the 
    panda example, and requires only one step.
    \item PGD \parencite{madry2018towards} is the standard strong baseline for evaluating defences.
\end{itemize}
\end{frame}

\begin{frame}
\frametitle{What the lab does}
\begin{itemize}
    % \item Model: ViT-Base, 16$\times$16 patches, ImageNet classes. Input: a COCO
    % photo of two cats, classified as ``Egyptian cat''. Target: ``daisy''.
    \item First unconstrained: optimise $\dd$ with Adam until the label flips.
    Then constrain it:
\end{itemize}
\begin{block}{Penalty plus box}
$$
\min_{\dd}\;
\Loss\big(\xx + \tilde\dd,\; y_\tgt\big) \;+\; \lambda\, \lVert \tilde\dd \rVert_2,
\qquad \tilde\dd = \clip_{[-\eps,\,\eps]}(\dd)
$$
\end{block}
\begin{itemize}
    \item $\lambda$ is a hyperparameter that controls trade-off between attack success and perturbation visibility.
\end{itemize}
\end{frame}

\begin{frame}
\frametitle{No gradients? Guess and check}
\begin{center}
\includegraphics[width=0.85\textwidth]{simba_gcv}
\figsrc{Google Cloud Vision before and after the attack. Figure from \textcite{guo2019simple}.}
\end{center}
\begin{block}{SimBA: only needs the returned probabilities}
Pick a random direction $\mathbf{q}$. Query $\xx + \dd \pm \eps\mathbf{q}$.
If the true class probability drops, keep the step. Repeat.
\end{block}
\begin{itemize}
    \item Google Cloud Vision: original top 3 labels removed in 70\% of 50
    images within 5\,000 queries.
    \item A closed API is no defence if it returns confidence scores.
\end{itemize}
\end{frame}

% \begin{frame}
% \frametitle{Surviving the camera: expectation over transformation}
% A perturbation tuned to one exact image breaks under a small crop, rotation
% or change in lighting. The physical attacks optimise over all of them at once.
% \begin{block}{Expectation over transformation (EOT)}
% $$
% \min_{\dd}\; \Ex_{t \sim T}\Big[\, \Loss\big(t(\xx + \dd),\; y_\tgt\big) \Big]
% $$
% $T$ is a distribution over transformations: viewpoint, distance, scale,
% lighting, printer colour error \parencite{athalye2018synthesizing}.
% \end{block}
% \begin{itemize}
%     \item In practice: sample a batch of transformations at every step and
%     average the loss. Still plain gradient descent.
%     \item The stop sign samples distances and angles, and masks the
%     perturbation into sticker shapes \parencite{eykholt2018robust}.
%     \item Watermarking is the same problem: an imperceptible mask
%     that must survive transformations applied to the image.
% \end{itemize}
% \end{frame}
